
\subsubsection{6.1 Key Finding: Label Correlation, Not Augmentation Invariance}

Our results are most parsimoniously explained by supervised label correlation as the dominant driver of spurious feature encoding, rather than augmentation-based instance-discriminability.

In ERM, the cross-entropy objective trains on data where background texture predicts bird species with 95\% reliability (Waterbirds). The model jointly encodes task and spurious features because both reduce training loss. MoCo-v3, being label-agnostic, encodes spurious features only insofar as they are instance-discriminative under the augmentation set. Background texture is instance-discriminative (different images have different backgrounds), but without label-correlation amplification, the resulting spurious encoding is measurably weaker (ratio = 1.027 vs 1.052, d = 5.68).

DINO's ERM-equivalent spurious encoding (d = 0.60, p\_bonf = 1.0) is consistent with the interpretation that DINO's momentum teacher produces class-correlated pseudo-targets that implicitly replicate supervised label selection pressure. Under this interpretation, it is the class-level signal --- whether explicit (ERM labels) or implicit (DINO teacher targets) --- that drives spurious feature encoding, not the augmentation scheme. This interpretation requires further verification (e.g., measuring mutual information between DINO teacher targets and spurious attribute labels) and remains a hypothesis rather than an established mechanism.

The augmentation ablation supports the complementary view that augmentation design can suppress spurious encoding when it explicitly removes the spurious attribute from the training signal. A 23.1\% ratio reduction (d = 32.4) when background is replaced per-view is consistent with contrastive objectives encoding spurious features to the degree they remain instance-discriminative after augmentation.

\subsubsection{6.2 Cross-Dataset Interaction}

The paradigm ranking reversal between Waterbirds and CelebA indicates that the relationship between pretraining objective and spurious feature encoding is not universal --- it depends on the type of spurious attribute. A plausible interpretation is that coarse texture attributes (background scenes) and demographic attributes (gender, distributed across face features) interact differently with each pretraining objective's encoding pressure. However, the current dataset sample (two datasets, two attribute types) is insufficient to characterize this interaction fully. The reversal is an empirical observation; the mechanism remains an open question.

\subsubsection{6.3 Limitations}

\textbf{L1: Original directional prediction refuted.} The initial hypothesis predicted that contrastive SSL would encode spurious features more strongly than ERM due to augmentation invariance. The data demonstrate the opposite with large effect size (d = 5.68 in the reversed direction). This is treated as a productive refutation; the empirical contributions are directionally opposite to the initial prediction but internally consistent.

\textbf{L2: No downstream worst-group accuracy evaluation.} The paper establishes paradigm differences in spurious/task probe accuracy ratio, but does not test whether lower spurious ratio translates to better worst-group accuracy when a linear head is trained via DFR-style fine-tuning \citep{Kirichenko2022}. Downstream consequences for practical robustness remain untested and should not be inferred from the probe results alone.

\textbf{L3: MoCo-v3 as contrastive representative.} SimCLR-v2 has no official PyTorch checkpoint; MoCo-v3 was used as the best available contrastive proxy. MoCo-v3 uses a momentum encoder and queue, which differ from SimCLR's direct two-view comparison. Results for contrastive SSL may not generalize to SimCLR specifically.

\textbf{L4: Scope.} Results apply to ResNet-50 backbones pretrained on ImageNet-1k, evaluated on Waterbirds and CelebA. ViT backbones, larger pretraining scales (e.g., CLIP), and additional spurious attribute types require separate study. The cross-dataset ranking reversal already indicates that scope matters.

\textbf{L5: CelebA test set size.} The Bonferroni-corrected CelebA pairwise test (h-e2) uses 720 balanced test images. A separate directional test (h-d1) uses a different balanced sample of 720 images; the two yield consistent qualitative conclusions (no significant ERM vs MoCo-v3 difference on CelebA), but the small balanced test size limits statistical power for detecting small effects.

\textbf{L6: Unverified citations.} Three citations --- Robinson et al. [2021], Wen et al. [2021], and Izmailov et al. [2022] --- have not been verified against their primary sources for title, venue, and claim correspondence. These are marked accordingly and must be verified before submission.

\subsubsection{6.4 Broader Implications}

These findings provide backbone selection guidance for fairness-critical applications: contrastive SSL backbones (particularly MoCo-v3 on texture-spurious datasets) partially suppress spurious features by design, without requiring group labels or robustness-specific interventions. This suppression is partial --- all paradigms produce ratios above 1.0, indicating that spurious attributes remain linearly decodable in all cases. The guidance is additive with, not a replacement for, existing robustness methods.

The augmentation ablation result (23.1\% ratio reduction) additionally suggests that augmentation design is an accessible lever for practitioners training SSL representations from scratch when the spurious attribute type is known in advance.

